\newif\ifglobal

%%%%%%%%%%%%%%%%%%%%%%%
%%% Compile to view %%%
%%%%%%%%%%%%%%%%%%%%%%%

%%%%%%%%%%%%%%%%%%%%%%%%%%%%%%%%%%%%%%%%%%%%%%%%%%%
%%% Readme:                                     %%%
%%% https://www.overleaf.com/read/bwdjvfjksvjy/ %%%
%%%%%%%%%%%%%%%%%%%%%%%%%%%%%%%%%%%%%%%%%%%%%%%%%%%

% >>> M?: marks a module setting number or important notice

% >>> M4: Set {./relative-path} to external files for this module <<<
\def \modulefiles {./23-DBGASSIST}
% To add an external file, use:
% (Replace '---' with actual filename)
% '\input{\modulefiles/tables/---.tex}' for tables
% '\includegraphics{\modulefiles/figures/---}' for figures
% '\subfile{\modulefiles/---}' for register files

% >>> M1: This module is in global project? <<<
% 'YES' - uncomment; 'NO' - comment out
\globaltrue

\ifglobal
    \documentclass[main\_\_CN.tex]{subfiles}

\else
    % Font "FandolSong-Regular" used by ctex does not have GSUB table.
    % It causes warnings that do not affect compilation and can be ignored.
    \PassOptionsToPackage{quiet}{fontspec}

    \documentclass[openany, 10pt]{book}

    % >>> M2: In ./chip-spec-settings/preamble-trm-module, also find
    %         settings for visibility of todo notes, watermark <<<
    \usepackage{preamble-trm-module}


    % Enable Chinese typesetting
    \usepackage{ctex}

    % >>> M3: Temporary labels in Chinese
    %         you can add your own labels there <<<
    \myexternaldocument{temp-labels-cn}

\fi %\ifglobal


\setcounter{secnumdepth}{4}


% >>> M5: If this module needs any updates in
% './chip-spec-settings/preamble-trm-module.sty'
% (include additional package, etc.), contact your module PM
% or create the file './chip-spec-settings/changelog.tex' make the
% necessary updates yourself and document them in this file <<<


\begin{document}


% Sets language into which pre-defined bits of text are translated
\selectlanguage{chinese}

% Do not delete this block
\ifglobal
% Add the following front matter if
% this module is outside of global project
\else
    \InputIfFileExists{readme}{}{}
    {\let\clearpage\relax \listoftodos}
    {\let\clearpage\relax \listofdonetodos}
    \tableofcontents
    %\listoftables
    %\listoffigures
\fi


%%%%%%%%%%%%%%%%%%%%%%%%%
%%% TRM Module Begins %%%
%%%%%%%%%%%%%%%%%%%%%%%%%

\chapter{辅助调试}
\label{mod:dbgassist}
\hypertarget{dbgassist}{}

\section{概述}

辅助调试模块提供一套调试功能，可用于软件开发时进行调试定位问题所在。

\section{主要特性}
\label{sec:dbgassist-features}

\begin{itemize}
    \item {\bfseries 读写监测}：监测一个高性能双核处理器（High-Performance CPU, HP CPU0 和 HP CPU1）总线是否在限定的存储器地址范围内进行读写操作，若在该地址范围内发生读写操作则触发中断。
    \item {\bfseries 栈指针 (SP) 监测}：监测栈指针是否超出限定的范围，若超出范围则产生中断。
    \item {\bfseries 程序计数器 (PC) 记录}：记录 PC，可以获得上一次 HP CPU0 或 HP CPU1 复位时的 PC 值。
    \item {\bfseries 总线访问记录}：记录总线访问信息，当 HP CPU0、HP CPU1 或 DMA 写了某个特殊值时，会记录此次写操作的总线类型、地址和 PC 值（仅记录 HP CPU0 或 HP CPU1 写操作的 PC），并将这些信息记录到 HP L2MEM 中。
\end{itemize}



\section{功能描述}
\label{sec:dbgassist-func-descr}

\subsection{区域读写监测}
为确认 HP CPU0 是否在某段地址范围（即区域）进行过读写行为，辅助调试模块能够监测 HP CPU0 总线，当总线在监测地址范围进行读写操作时会触发中断，详情请见章节 \ref{mod:riscvcpu} \textit{\nameref{mod:riscvcpu}} 中的表 \ref{tab:riscvcpu-mem-blocks}。 当 HP CPU0 总线访问数据空间时，本章节将其称为数据总线；而当 HP CPU0 总线访问 AHB 外设空间时，本章节将其称为外设总线。辅助调试模块可同时监测数据总线在两段地址范围内的读写行为，假设命名为 HP CPU0 数据总线区域 0 和 HP CPU0 数据总线区域 1，外设总线也可同时监测两段地址范围，假设命名为 HP CPU0 外设总线区域 0 和 HP CPU0 外设总线区域 1。区域 0 和区域 1 根据需要进行配置。

HP CPU1 同理。

\subsection{栈指针监测}
为防止 HP CPU0 栈溢出或者错误的压栈出栈，辅助调试模块能够监测栈指针。当 HP CPU0 栈指针超过 HP CPU0 栈指针监测区域限定的上下边界时，模块会记录 PC 指针并产生中断，边界由软件进行配置。

HP CPU1 同理。

\subsection{PC 记录}
在某些时候软件开发者希望知道上次 HP CPU0 复位时的 PC 指针。比如，在程序卡死只能复位时，开发者可能希望读取复位时的 PC 指针以便知道程序在哪里卡死，然后进行调试。辅助调试模块可以记录 HP CPU0 复位时的 PC，方便开发者进行调试。

HP CPU1 同理。

\subsection{CPU/DMA 总线访问记录}
辅助调试模块能够实时记录 HP CPU0 数据总线、 HP CPU1 数据总线和 DMA 总线的写行为，当总线在某个范围内发生写操作或者在某个范围内写某个特定的值时，此次操作的总线类型、地址和 PC（仅记录 HP CPU0 或 HP CPU1 写操作的 PC）等信息会被记录下来，并按照一定的格式存储到 HP L2MEM 中。此处提到的地址范围只能处于 HP SPM 或 HP L2MEM 的范围区间。


%%%%%%%%%%%%%%%%%%
%%% Interrupts %%%
%%%%%%%%%%%%%%%%%%

\section{中断}
\label{sec:dbgassist-interrupts}


以下辅助调试模块的中断源会生成 ASSIST\_DEBUG\_INT 中断信号。

\begin{itemize}
\item \label{int:ASSISTDEBUGCORE0AREADRAM00RDINT}ASSIST\_DEBUG\_CORE\_0\_AREA\_DRAM0\_0\_RD\_INT：当 HP CPU0 数据总线在 HP CPU0 数据总线区域 0 读操作时触发。
\item \label{int:ASSISTDEBUGCORE0AREADRAM00WRINT}ASSIST\_DEBUG\_CORE\_0\_AREA\_DRAM0\_0\_WR\_INT：当 HP CPU0 数据总线在 HP CPU0 数据总线区域 0 写操作时触发。
\item \label{int:ASSISTDEBUGCORE0AREADRAM01RDINT}ASSIST\_DEBUG\_CORE\_0\_AREA\_DRAM0\_1\_RD\_INT：当 HP CPU0 数据总线在 HP CPU0 数据总线区域 1 读操作时触发。
\item \label{int:ASSISTDEBUGCORE0AREADRAM01WRINT}ASSIST\_DEBUG\_CORE\_0\_AREA\_DRAM0\_1\_WR\_INT：当 HP CPU0 数据总线在 HP CPU0 数据总线区域 1 写操作时时触发。
\item \label{int:ASSISTDEBUGCORE0AREAPIF0RDINT}ASSIST\_DEBUG\_CORE\_0\_AREA\_PIF\_0\_RD\_INT：当 HP CPU0 外设总线在 HP CPU0 外设总线区域 0 读操作时触发。
\item \label{int:ASSISTDEBUGCORE0AREAPIF0WRINT}ASSIST\_DEBUG\_CORE\_0\_AREA\_PIF\_0\_WR\_INT：当 HP CPU0 外设总线在 HP CPU0 外设总线区域 0 写操作时时触发。
\item \label{int:ASSISTDEBUGCORE0AREAPIF1RDINT}ASSIST\_DEBUG\_CORE\_0\_AREA\_PIF\_1\_RD\_INT：当 HP CPU0 外设总线在 HP CPU0 外设总线区域 1 读操作时触发。
\item \label{int:ASSISTDEBUGCORE0AREAPIF1WRINT}ASSIST\_DEBUG\_CORE\_0\_AREA\_PIF\_1\_WR\_INT：当 HP CPU0 外设总线在 HP CPU0 外设总线区域 1 写操作时时触发。
\item
\label{int:ASSISTDEBUGCORE0SPSPILLMININT}ASSIST\_DEBUG\_CORE\_0\_SP\_SPILL\_MIN\_INT：当 HP CPU0 栈指针小于 HP CPU0 栈指针监测区域的下边界时触发。
\item
\label{int:ASSISTDEBUGCORE0SPSPILLMAXINT}ASSIST\_DEBUG\_CORE\_0\_SP\_SPILL\_MAX\_INT：当 HP CPU0 栈指针大于 HP CPU0 栈指针监测区域的上边界时触发。
\item \label{int:ASSISTDEBUGCORE1AREADRAM00RDINT}ASSIST\_DEBUG\_CORE\_1\_AREA\_DRAM0\_0\_RD\_INT：当 HP CPU1 数据总线在 HP CPU1 数据总线区域 0 读操作时触发。
\item \label{int:ASSISTDEBUGCORE1AREADRAM00WRINT}ASSIST\_DEBUG\_CORE\_1\_AREA\_DRAM0\_0\_WR\_INT：当 HP CPU1 数据总线在 HP CPU1 数据总线区域 0 写操作时触发。
\item \label{int:ASSISTDEBUGCORE1AREADRAM01RDINT}ASSIST\_DEBUG\_CORE\_1\_AREA\_DRAM0\_1\_RD\_INT：当 HP CPU1 数据总线在 HP CPU1 数据总线区域 1 读操作时触发。
\item \label{int:ASSISTDEBUGCORE1AREADRAM01WRINT}ASSIST\_DEBUG\_CORE\_1\_AREA\_DRAM0\_1\_WR\_INT：当 HP CPU1 数据总线在 HP CPU1 数据总线区域 1 写操作时时触发。
\item \label{int:ASSISTDEBUGCORE1AREAPIF0RDINT}ASSIST\_DEBUG\_CORE\_1\_AREA\_PIF\_0\_RD\_INT：当 HP CPU1 外设总线在 HP CPU1 外设总线区域 0 读操作时触发。
\item \label{int:ASSISTDEBUGCORE1AREAPIF0WRINT}ASSIST\_DEBUG\_CORE\_1\_AREA\_PIF\_0\_WR\_INT：当 HP CPU1 外设总线在 HP CPU1 外设总线区域 0 写操作时时触发。
\item \label{int:ASSISTDEBUGCORE1AREAPIF1RDINT}ASSIST\_DEBUG\_CORE\_1\_AREA\_PIF\_1\_RD\_INT：当 HP CPU1 外设总线在 HP CPU1 外设总线区域 1 读操作时触发。
\item \label{int:ASSISTDEBUGCORE1AREAPIF1WRINT}ASSIST\_DEBUG\_CORE\_1\_AREA\_PIF\_1\_WR\_INT：当 HP CPU1 外设总线在 HP CPU1 外设总线区域 1 写操作时时触发。
\item
\label{int:ASSISTDEBUGCORE1SPSPILLMININT}ASSIST\_DEBUG\_CORE\_1\_SP\_SPILL\_MIN\_INT：当 HP CPU1 栈指针小于 HP CPU1 栈指针监测区域的下边界时触发。
\item
\label{int:ASSISTDEBUGCORE1SPSPILLMAXINT}ASSIST\_DEBUG\_CORE\_1\_SP\_SPILL\_MAX\_INT：当 HP CPU1 栈指针大于 HP CPU1 栈指针监测区域的上边界时触发。
\end{itemize}

\begin{tiplisting}
对于中断、中断信号、中断源 及其关系的详细描述，可以查看章节 \ref{mod:intmtrx} \textit{\nameref{mod:intmtrx}} > Section \ref{sec:intmtx-int-terminology} \textit{\nameref{sec:intmtx-int-terminology}}。
\end{tiplisting}

每个中断源都有一组通用的配置寄存器，详见章节 \hyperref[interrupt-config-registers]{\textit{中断配置寄存器}}。具体寄存器名称请查看 \ref{sec:dbgassist-reg-summ} \textit{\nameref{sec:dbgassist-reg-summ}}。


\section{工作流程}
\label{sec:dbgassist-recommended-operation}

\subsection{区域监测和栈指针监测配置}

区域监测可监测 HP CPU0 和 HP CPU1 的数据总线和外设总线的读写行为，每个总线可同时监测两个区域。详细监测模式如下。

\begin{itemize}
    \item 数据总线读写监测
        \begin{itemize}
            \item 监测 HP CPU0 数据总线在 HP CPU0 数据总线区域 0 是否进行读操作
            \item 监测 HP CPU0 数据总线在 HP CPU0 数据总线区域 0 是否进行写操作
            \item 监测 HP CPU0 数据总线在 HP CPU0 数据总线区域 1 是否进行读操作
            \item 监测 HP CPU0 数据总线在 HP CPU0 数据总线区域 1 是否进行写操作
            \item 监测 HP CPU1 数据总线在 HP CPU1 数据总线区域 0 是否进行读操作
            \item 监测 HP CPU1 数据总线在 HP CPU1 数据总线区域 0 是否进行写操作
            \item 监测 HP CPU1 数据总线在 HP CPU1 数据总线区域 1 是否进行读操作
            \item 监测 HP CPU1 数据总线在 HP CPU1 数据总线区域 1 是否进行写操作
        \end{itemize}
    \item 外设总线读写监测
        \begin{itemize}
            \item 监测 HP CPU0 外设总线在 HP CPU0 外设总线区域 0 是否进行读操作
            \item 监测 HP CPU0 外设总线在 HP CPU0 外设总线区域 0 是否进行写操作
            \item 监测 HP CPU0 外设总线在 HP CPU0 外设总线区域 1 是否进行读操作
            \item 监测 HP CPU0 外设总线在 HP CPU0 外设总线区域 1 是否进行写操作
            \item 监测 HP CPU1 外设总线在 HP CPU1 外设总线区域 0 是否进行读操作
            \item 监测 HP CPU1 外设总线在 HP CPU1 外设总线区域 0 是否进行写操作
            \item 监测 HP CPU1 外设总线在 HP CPU1 外设总线区域 1 是否进行读操作
            \item 监测 HP CPU1 外设总线在 HP CPU1 外设总线区域 1 是否进行写操作
        \end{itemize}
    \item 栈指针监测
        \begin{itemize}
            \item 监测 HP CPU0 栈指针是否大于 HP CPU0 栈指针监测区域上边界
            \item 监测 HP CPU0 栈指针是否小于 HP CPU0 栈指针监测区域下边界
            \item 监测 HP CPU1 栈指针是否大于 HP CPU1 栈指针监测区域上边界
            \item 监测 HP CPU1 栈指针是否小于 HP CPU1 栈指针监测区域下边界
        \end{itemize}
\end{itemize}

区域监测和栈指针监测的配置流程如下：
\begin{enumerate}
    \item 配置监测区域和栈指针
        \begin{itemize}
            \item HP CPU0 数据总线区域 0 由 \hyperref[regdesc:ASSISTDEBUGCORE0AREADRAM00MINREG]{ASSIST\_DEBUG\_CORE\_0\_AREA\_DRAM0\_0\_MIN\_REG} 和\\ \hyperref[regdesc:ASSISTDEBUGCORE0AREADRAM00MAXREG]{ASSIST\_DEBUG\_CORE\_0\_AREA\_DRAM0\_0\_MAX\_REG} 进行配置。
            \item HP CPU0 数据总线区域 1 由 \hyperref[regdesc:ASSISTDEBUGCORE0AREADRAM01MINREG]{ASSIST\_DEBUG\_CORE\_0\_AREA\_DRAM0\_1\_MIN\_REG} 和\\ \hyperref[regdesc:ASSISTDEBUGCORE0AREADRAM01MAXREG]{ASSIST\_DEBUG\_CORE\_0\_AREA\_DRAM0\_1\_MAX\_REG} 进行配置。
            \item HP CPU0 外设总线区域 0 由 \hyperref[regdesc:ASSISTDEBUGCORE0AREAPIF0MINREG]{ASSIST\_DEBUG\_CORE\_0\_AREA\_PIF\_0\_MIN\_REG} 和\\ \hyperref[regdesc:ASSISTDEBUGCORE0AREAPIF0MAXREG]{ASSIST\_DEBUG\_CORE\_0\_AREA\_PIF\_0\_MAX\_REG} 进行配置。
            \item HP CPU0 外设总线区域 1 由 \hyperref[regdesc:ASSISTDEBUGCORE0AREAPIF1MINREG]{ASSIST\_DEBUG\_CORE\_0\_AREA\_PIF\_1\_MIN\_REG} 和\\ \hyperref[regdesc:ASSISTDEBUGCORE0AREAPIF1MAXREG]{ASSIST\_DEBUG\_CORE\_0\_AREA\_PIF\_1\_MAX\_REG} 进行配置。
            \item HP CPU0 栈指针监测区域由 \hyperref[regdesc:ASSISTDEBUGCORE0SPMINREG]{ASSIST\_DEBUG\_CORE\_0\_SP\_MIN\_REG} 和\\ \hyperref[regdesc:ASSISTDEBUGCORE0SPMAXREG]{ASSIST\_DEBUG\_CORE\_0\_SP\_MAX\_REG} 进行配置。
            \item HP CPU1 数据总线区域 0 由 \hyperref[regdesc:ASSISTDEBUGCORE1AREADRAM00MINREG]{ASSIST\_DEBUG\_CORE\_1\_AREA\_DRAM0\_0\_MIN\_REG} 和\\ \hyperref[regdesc:ASSISTDEBUGCORE1AREADRAM00MAXREG]{ASSIST\_DEBUG\_CORE\_1\_AREA\_DRAM0\_0\_MAX\_REG} 进行配置。
            \item HP CPU1 数据总线区域 1 由 \hyperref[regdesc:ASSISTDEBUGCORE1AREADRAM01MINREG]{ASSIST\_DEBUG\_CORE\_1\_AREA\_DRAM0\_1\_MIN\_REG} 和\\ \hyperref[regdesc:ASSISTDEBUGCORE1AREADRAM01MAXREG]{ASSIST\_DEBUG\_CORE\_1\_AREA\_DRAM0\_1\_MAX\_REG} 进行配置。
            \item HP CPU1 外设总线区域 0 由 \hyperref[regdesc:ASSISTDEBUGCORE1AREAPIF0MINREG]{ASSIST\_DEBUG\_CORE\_1\_AREA\_PIF\_0\_MIN\_REG} 和\\ \hyperref[regdesc:ASSISTDEBUGCORE1AREAPIF0MAXREG]{ASSIST\_DEBUG\_CORE\_1\_AREA\_PIF\_0\_MAX\_REG} 进行配置。
            \item HP CPU1 外设总线区域 1 由 \hyperref[regdesc:ASSISTDEBUGCORE1AREAPIF1MINREG]{ASSIST\_DEBUG\_CORE\_1\_AREA\_PIF\_1\_MIN\_REG} 和\\ \hyperref[regdesc:ASSISTDEBUGCORE1AREAPIF1MAXREG]{ASSIST\_DEBUG\_CORE\_1\_AREA\_PIF\_1\_MAX\_REG} 进行配置。
            \item HP CPU1 栈指针监测区域由 \hyperref[regdesc:ASSISTDEBUGCORE1SPMINREG]{ASSIST\_DEBUG\_CORE\_1\_SP\_MIN\_REG} 和\\ \hyperref[regdesc:ASSISTDEBUGCORE1SPMAXREG]{ASSIST\_DEBUG\_CORE\_1\_SP\_MAX\_REG} 进行配置。
        \end{itemize}
    \item 配置中断
    \begin{itemize}
        \item 配置 \hyperref[regdesc:ASSISTDEBUGCORE0INTRENAREG]{ASSIST\_DEBUG\_CORE\_0\_INTR\_ENA\_REG} 或 \hyperref[regdesc:ASSISTDEBUGCORE1INTRENAREG]{ASSIST\_DEBUG\_CORE\_1\_INTR\_ENA\_REG} 用于使能不同模式的中断。
        \item 读取 \hyperref[regdesc:ASSISTDEBUGCORE0INTRRAWREG]{ASSIST\_DEBUG\_CORE\_0\_INTR\_RAW\_REG} 或 \hyperref[regdesc:ASSISTDEBUGCORE1INTRRAWREG]{ASSIST\_DEBUG\_CORE\_1\_INTR\_RAW\_REG} 用于查询不同模式的原始中断状态。
        \item 配置 \hyperref[regdesc:ASSISTDEBUGCORE0INTRCLRREG]{ASSIST\_DEBUG\_CORE\_0\_INTR\_CLR\_REG} 或 \hyperref[regdesc:ASSISTDEBUGCORE1INTRCLRREG]{ASSIST\_DEBUG\_CORE\_1\_INTR\_CLR\_REG} 用于清除不同模式的中断。
        \end{itemize}
    \item 配置 \hyperref[regdesc:ASSISTDEBUGCORE0MONTRENAREG]{ASSIST\_DEBUG\_CORE\_0\_MONTR\_ENA\_REG} 或 \hyperref[regdesc:ASSISTDEBUGCORE1MONTRENAREG]{ASSIST\_DEBUG\_CORE\_1\_MONTR\_ENA\_REG} 使能不同的监测模式，可同时使能。
\end{enumerate}

 比如，若想监测 HP CPU0 数据总线地址 A 到地址 B 范围内是否进行过写操作，可通过 HP CPU0 数据总线区域 0 或者区域 1 进行监测，以 HP CPU0 数据总线区域 0 为例：
\begin{enumerate}
    \item 配置 \hyperref[fielddesc:ASSISTDEBUGCORE0RCDPDEBUGEN]{ASSIST\_DEBUG\_CORE\_0\_RCD\_PDEBUGEN} 为 1 ，使能 HP CPU0 更新 PC 信号给辅助调试模块。
    \item 配置 \hyperref[regdesc:ASSISTDEBUGCORE0AREADRAM00MINREG]{ASSIST\_DEBUG\_CORE\_0\_AREA\_DRAM0\_0\_MIN\_REG} 为 A 地址。
    \item 配置 \hyperref[regdesc:ASSISTDEBUGCORE0AREADRAM00MAXREG]{ASSIST\_DEBUG\_CORE\_0\_AREA\_DRAM0\_0\_MAX\_REG} 为 B 地址。
    \item 配置 \hyperref[regdesc:ASSISTDEBUGCORE0INTRENAREG]{ASSIST\_DEBUG\_CORE\_0\_INTR\_ENA\_REG} bit[1] 使能 HP CPU0 数据总线区域 0 写监测中断。
    \item 配置 \hyperref[regdesc:ASSISTDEBUGCORE0MONTRENAREG]{ASSIST\_DEBUG\_CORE\_0\_MONTR\_ENA\_REG} bit[1] 使能 HP CPU0 数据总线区域 0 写监测。
    \item 配置中断矩阵将 ASSIST\_DEBUG\_INTR 映射到 HP CPU0 中断或 HP CPU1 中断（参考章节 \ref{mod:intmtrx} \textit{\nameref{mod:intmtrx}})）。
    \item 发生中断后
    \begin{itemize}
        \item 读取 \hyperref[regdesc:ASSISTDEBUGCORE0INTRRAWREG]{ASSIST\_DEBUG\_CORE\_0\_INTR\_RAW\_REG} 和 \hyperref[regdesc:ASSISTDEBUGCORE1INTRRAWREG]{ASSIST\_DEBUG\_CORE\_1\_INTR\_RAW\_REG} 确认是什么操作触发的中断。
        \item 如果是 HP CPU0 区域监测引发的中断，读取 \hyperref[fielddesc:ASSISTDEBUGCORE0AREAPC]{ASSIST\_DEBUG\_CORE\_0\_AREA\_PC} 可获取触发中断时刻的 HP CPU0 PC 值，读取 \hyperref[fielddesc:ASSISTDEBUGCORE0AREASP]{ASSIST\_DEBUG\_CORE\_0\_AREA\_SP} 可获取触发中断时刻的 HP CPU0 SP 值。
        \item 如果是栈指针监测引发的中断，读取 \hyperref[fielddesc:ASSISTDEBUGCORE0SPPC]{ASSIST\_DEBUG\_CORE\_0\_SP\_PC} 可获取触发中断时刻的 HP CPU0 PC 值。
        \item 配置 \hyperref[regdesc:ASSISTDEBUGCORE0INTRCLRREG]{ASSIST\_DEBUG\_CORE\_0\_INTR\_CLR\_REG} 和 \hyperref[regdesc:ASSISTDEBUGCORE1INTRCLRREG]{ASSIST\_DEBUG\_CORE\_1\_INTR\_CLR\_REG} 不同位可清除不同的中断。
    \end{itemize}

\end{enumerate}


\subsection{PC 记录配置}

配置 \hyperref[fielddesc:ASSISTDEBUGCORE0RCDPDEBUGEN]{ASSIST\_DEBUG\_CORE\_0\_RCD\_PDEBUGEN} 为 1 ，使能 HP CPU0 更新 PC 信号给辅助调试模块。同时当 \hyperref[fielddesc:ASSISTDEBUGCORE0RCDRECORDEN]{ASSIST\_DEBUG\_CORE\_0\_RCD\_RECORDEN} 配置为 1 时，\\\hyperref[regdesc:ASSISTDEBUGCORE0RCDPDEBUGPCREG]{ASSIST\_DEBUG\_CORE\_0\_RCD\_PDEBUGPC\_REG} 会去记录 HP CPU0 PC 信号，同时 \hyperref[regdesc:ASSISTDEBUGCORE0RCDPDEBUGSPREG]{ASSIST\_DEBUG\_CORE\_0\_RCD\_PDEBUGSP\_REG} 会记录 HP CPU0 SP 值，否则保持原值。

当 HP CPU0 发生复位时，\hyperref[regdesc:ASSISTDEBUGCORE0RCDENREG]{ASSIST\_DEBUG\_CORE\_0\_RCD\_EN\_REG} 会被复位，但是\\ \hyperref[regdesc:ASSISTDEBUGCORE0RCDPDEBUGPCREG]{ASSIST\_DEBUG\_CORE\_0\_RCD\_PDEBUGPC\_REG} 和 \hyperref[regdesc:ASSISTDEBUGCORE0RCDPDEBUGSPREG]{ASSIST\_DEBUG\_CORE\_0\_RCD\_PDEBUGSP\_REG} 不会被复位，因此这两个寄存器会一直保持 HP CPU0 复位时刻的 PC 值 和 SP 值。

HP CPU1 同理。


\subsection{CPU/DMA 总线访问记录配置}

\subsubsection{HP CPU0/1 总线访问记录配置}
\label{sec:hp_cpu_access_mon}

HP CPU0/1 总线访问记录（SPM 监控功能）的配置流程如下：

\begin{enumerate}
    \item 配置监测地址范围
        \begin{itemize}
            \item 配置 \hyperref[regdesc:SPMMEMMONITORLOGMINREG]{SPM\_MEM\_MONITOR\_LOG\_MIN\_REG} 和 \hyperref[regdesc:SPMMEMMONITORLOGMAXREG]{SPM\_MEM\_MONITOR\_LOG\_MAX\_REG} 确定地址范围。监测地址范围需处于 0x3010\_0000 \verb+~+ 0x3010\_1FFF 区间内。
        \end{itemize}
    \item 配置监测模式 (\hyperref[fielddesc:SPMMEMMONITORLOGMODE]{SPM\_MEM\_MONITOR\_LOG\_MODE})
       \begin{itemize}
            \item 写监测（监测总线是否发生写操作）
           \item 字 (word) 监测（监测总线是否写某个特定 word）
           \item 半字 (halfword) 监测（监测总线是否写某个特定 halfword）
           \item 字节 (byte) 监测（监测总线是否写某个特定 byte）
       \end{itemize}
    \item 配置需要监测的特殊值
        \begin{itemize}
            \item 监测模式为 word 监测时，\hyperref[regdesc:SPMMEMMONITORLOGCHECKDATAREG]{SPM\_MEM\_MONITOR\_LOG\_CHECK\_DATA\_REG} 即为监测的特定 word。
            \item 监测模式为 halfword 监测时，\hyperref[regdesc:SPMMEMMONITORLOGCHECKDATAREG]{SPM\_MEM\_MONITOR\_LOG\_CHECK\_DATA\_REG} 的 [15:0] 即为监测的特定 halfword。
            \item 监测模式为 byte 监测时，\hyperref[regdesc:SPMMEMMONITORLOGCHECKDATAREG]{SPM\_MEM\_MONITOR\_LOG\_CHECK\_DATA\_REG} 的 [7:0] 即为监测的特定 byte。
            \item \hyperref[regdesc:SPMMEMMONITORLOGDATAMASKREG]{SPM\_MEM\_MONITOR\_LOG\_DATA\_MASK\_REG} 用于屏蔽 \hyperref[regdesc:SPMMEMMONITORLOGCHECKDATAREG]{SPM\_MEM\_MONITOR\_LOG\_CHECK\_DATA\_REG} 的相应 byte。当某一 byte 被屏蔽，表示该 byte 可为任意的值，比如 word 监测，\hyperref[regdesc:SPMMEMMONITORLOGCHECKDATAREG]{SPM\_MEM\_MONITOR\_LOG\_CHECK\_DATA\_REG} 配置为 0x01020304，\hyperref[regdesc:SPMMEMMONITORLOGDATAMASKREG]{SPM\_MEM\_MONITOR\_LOG\_DATA\_MASK\_REG} 配置为 0x1，因此只要总线写 0x010203XX 都会被记录。
        \end{itemize}
    \item 配置记录信息的存储地址范围
        \begin{itemize}
            \item \hyperref[regdesc:SPMMEMMONITORLOGMEMSTARTREG]{SPM\_MEM\_MONITOR\_LOG\_MEM\_START\_REG} 和 \hyperref[regdesc:SPMMEMMONITORLOGMEMENDREG]{SPM\_MEM\_MONITOR\_LOG\_MEM\_END\_REG} 为存储地址范围配置寄存器。 记录信息的存储地址范围为 0x4FF0\_0000 \verb+~+ 0x4FFB\_FFFF。
            \item 置位 \hyperref[regdesc:SPMMEMMONITORLOGMEMADDRUPDATEREG]{SPM\_MEM\_MONITOR\_LOG\_MEM\_ADDR\_UPDATE\_REG}，使 \hyperref[regdesc:SPMMEMMONITORLOGMEMCURRENTADDRREG]{SPM\_MEM\_MONITOR\_LOG\_MEM\_CURRENT\_ADDR\_REG} 更新为 \hyperref[regdesc:SPMMEMMONITORLOGMEMSTARTREG]{SPM\_MEM\_MONITOR\_LOG\_MEM\_START\_REG}。
            \item 配置辅助调试模块的 SPM 监控功能对 HP L2MEM 的使用权限，只有开启了辅助调试模块对 HP L2MEM 的使用权限才能访问 HP L2MEM。关于如何配置，详情见章节 \ref{mod:permctrl} \textit{\nameref{mod:permctrl}}) 的 \hyperref[regdesc:TEEDMASPMMONPMSWREG]{TEE\_DMA\_SPM\_MON\_PMS\_W\_REG} 及相关内容。
        \end{itemize}
    \item \label{other:dbgassist-item-log-config-step5} 配置存储记录信息的内存的存储模式：loop 模式和非 loop 模式
        \begin{itemize}
            \item loop 模式下，将记录信息循环写进配置地址范围内，当写到结束地址时回到起始地址开始写，覆盖之前记录的信息。置位 \hyperref[fielddesc:SPMMEMMONITORLOGMEMLOOPENABLE]{SPM\_MEM\_MONITOR\_LOG\_MEM\_LOOP\_ENABLE} 使用 loop 模式。\\
            例如，存储记录信息的地址范围为 0 \verb+~+ 4 ，总线访问过程中有 1 \verb+~+ 10 共十次写记录信息的操作，那么第 5 次写操作写到地址 4 后，第 6 次写操作会回到地址 0 开始写，第 6 \verb+~+ 10 写操作会覆盖之前的第 1 \verb+~+ 5 写操作的数据。
            \item 非 loop 模式下，当将记录信息写到结束地址后，会一直停留在结束地址，并不再写记录信息，不会覆盖最开始的记录信息，并且丢弃后面的记录信息。清零 \hyperref[fielddesc:SPMMEMMONITORLOGMEMLOOPENABLE]{SPM\_MEM\_MONITOR\_LOG\_MEM\_LOOP\_ENABLE} 使用非 loop 模式。\\
            例如，存储记录信息的地址范围为 0 \verb+~+ 4 ，总线访问过程中有 1 \verb+~+ 10 共十次写记录信息的操作，那么第 5 次写操作写到地址 4 后，第 6 \verb+~+ 10 次写操作的地址会一直在地址 4 上，但不再进行写操作，因此地址范围 0 \verb+~+ 4 保留前 5 次写操作信息，并且丢弃第 6 \verb+~+ 10 次写操作信息。
        \end{itemize}
     \item 配置总线使能
        \begin{itemize}
            \item 配置 \hyperref[fielddesc:SPMMEMMONITORLOGCOREENA]{SPM\_MEM\_MONITOR\_LOG\_CORE\_ENA} 使能 HP CPU0 或 HP CPU1 总线访问记录，可同时使能。
        \end{itemize}
\end{enumerate}

辅助调试模块会先将 HP CPU0 或 HP CPU1 的记录信息缓存到一个内部 buffer A 里，然后从 buffer 取出记录信息再存储到配置的存储区域内。当监测行为被连续触发时，连续产生的记录数据包可能会使得 buffer 满而不能缓存新的记录数据包，此时模块会舍弃掉不能及时缓存到 buffer 的记录数据包，并在 buffer 不满的时候缓存一个 HP CPU LOST 数据包作为替代。然而仅能知道 HP CPU LOST 数据包生成时所舍弃的记录数据包的总线类型，不能知道在此之前舍弃了多少个记录数据包。

当总线访问记录结束后需要从内存中读取记录数据并进行解码。记录数据里只有三种包格式，HP CPU0 数据包、 HP CPU1 数据包和 HP CPU LOST 数据包，HP CPU0/1 数据包分别对应 HP CPU0/1 数据总线记录信息，三种包格式如表 \ref{tab:dbgassist-hpcpu0-package} 和 \ref{tab:dbgassist-hpcpu-lost-package} 所示：

    \begin{longtable}[c]{|L{3cm}|L{3cm}|L{3cm}|L{3cm}|l|}
    \caption{HP CPU0/1 包格式}\label{tab:dbgassist-hpcpu0-package}
    \\\hline
    \rowcolor{lightgray}
    {\bfseries Bit[63:33]} & {\bfseries Bit[32]} & {\bfseries Bit[31:3]} & {\bfseries Bit[2:1]} & {\bfseries Bit[0]} \\\hline
    pc\_offset & anchored(1) & addr\_offset & format & anchored(0) \\\hline
    \end{longtable}

    \begin{longtable}[c]{|L{3cm}|L{3cm}|L{3cm}|l|}
    \caption{LOST 包格式}\label{tab:dbgassist-hpcpu-lost-package}
    \\\hline
    \rowcolor{lightgray}
    {\bfseries Bit[31:7]} & {\bfseries Bit[6:3]} & {\bfseries Bit[2:1]} & {\bfseries Bit[0]} \\\hline
    reserved & lost\_source & format & anchored(0) \\\hline
    \end{longtable}

从包格式可以看出，HP CPU0/1 数据包共 64 位， HP CPU LOST 数据包共 32 位。下面介绍各个域的含义：
\begin{itemize}
    \item {\bfseries format} 表示此次数据包的类型，0 表示 HP CPU0 数据包，1 表示 HP CPU1 数据包，2 表示 HP CPU LOST 数据包，3 为保留值。
    \item {\bfseries pc\_offset} 记录了 HP CPU0/1 的 PC 指针的偏移量，实际 PC = pc\_offset + 0x0000\_0000。
    \item {\bfseries addr\_offset} 记录了此次写操作的地址偏移，实际地址 = addr\_offset + \{\hyperref[regdesc:SPMMEMMONITORLOGMINREG]{SPM\_MEM\_MONITOR\_LOG\_MIN\_REG}[31:4], 4'h0\}。
    \item {\bfseries lost\_source} 记录了 HP CPU LOST 数据包生成时舍弃了哪些数据包：
    \begin{itemize}
        \item Bit[4:3]：0 表示未舍弃 HP CPU0 数据包，3 表示舍弃了 HP CPU0 数据包，其他值为保留值。
        \item Bit[6:5]：0 表示未舍弃 HP CPU1 数据包，3 表示舍弃了 HP CPU1 数据包，其他值为保留值。
    \end{itemize}
    \item {\bfseries anchored} 表示该 32 位数据在数据包中的位置，0 表示该 32 位数据是数据包的低 32 位，1 表示该 32 位数据是数据包的高 32 位。

\end{itemize}
模块内部 buffer 的数据宽度是 32 位。当同时使能 HP CPU0 和 HP CPU1 总线访问记录并且在某一时刻同时生成了记录数据，先将 HP CPU0 数据包缓存到 buffer 中，再缓存 HP CPU1 数据包。辅助调试模块会自动地将缓存数据从 buffer 取出，并将取出的数据再以 32 位数据宽度存入指定范围的 memory 里。
当采用 loop 模式时，若在配置的存储地址范围内循环了若干次，可能存在残留数据干扰解析，即 HP CPU0/1 数据包的低 32 位数据被覆盖而使得高 32 位数据成为残留数据，因此确定第一个有效数据包的位置时必须过滤可能存在的残留数据。使用 \hyperref[regdesc:SPMMEMMONITORLOGMEMCURRENTADDRREG]{SPM\_MEM\_MONITOR\_LOG\_MEM\_CURRENT\_ADDR\_REG} 确定了数据包的起始位置之后，检查数据包 anchored 位的数值，是 0 则保留，是 1 则舍弃。

数据包解析流程如下：

\begin{itemize}
    \item \hyperref[fielddesc:SPMMEMMONITORLOGMEMFULLFLAG]{SPM\_MEM\_MONITOR\_LOG\_MEM\_FULL\_FLAG} 确定数据是否溢出配置范围。
    \begin{itemize}
        \item 如果未溢出，则数据读取的地址范围为 \hyperref[regdesc:SPMMEMMONITORLOGMEMSTARTREG]{SPM\_MEM\_MONITOR\_LOG\_MEM\_START\_REG}  \verb+~+ \hyperref[regdesc:SPMMEMMONITORLOGMEMCURRENTADDRREG]{SPM\_MEM\_MONITOR\_LOG\_MEM\_CURRENT\_ADDR\_REG} - 4 ；
        \item 若溢出并且开启了 loop 模式，则数据读取的地址范围为 \hyperref[regdesc:SPMMEMMONITORLOGMEMCURRENTADDRREG]{SPM\_MEM\_MONITOR\_LOG\_MEM\_CURRENT\_ADDR\_REG} \verb+~+ \hyperref[regdesc:SPMMEMMONITORLOGMEMENDREG]{SPM\_MEM\_MONITOR\_LOG\_MEM\_END\_REG} 和 \hyperref[regdesc:SPMMEMMONITORLOGMEMSTARTREG]{SPM\_MEM\_MONITOR\_LOG\_MEM\_START\_REG} \verb+~+ \hyperref[regdesc:SPMMEMMONITORLOGMEMCURRENTADDRREG]{SPM\_MEM\_MONITOR\_LOG\_MEM\_CURRENT\_ADDR\_REG} - 4；
        \item 若溢出并且未开启 loop 模式，则数据读取的地址范围为 \hyperref[regdesc:SPMMEMMONITORLOGMEMSTARTREG]{SPM\_MEM\_MONITOR\_LOG\_MEM\_START\_REG} \verb+~+ \hyperref[regdesc:SPMMEMMONITORLOGMEMENDREG]{SPM\_MEM\_MONITOR\_LOG\_MEM\_END\_REG}。
    \end{itemize}
    \item 从起始地址处开始读取数据并解析，每次读取 32 位。

\end{itemize}

数据解析完成之后需要通过配置 \hyperref[fielddesc:SPMMEMMONITORCLRLOGMEMFULLFLAG]{SPM\_MEM\_MONITOR\_CLR\_LOG\_MEM\_FULL\_FLAG} 清除  \hyperref[fielddesc:SPMMEMMONITORLOGMEMFULLFLAG]{SPM\_MEM\_MONITOR\_LOG\_MEM\_FULL\_FLAG} 标志位。



\subsubsection{DMA 总线访问记录配置}

辅助调试模块共支持 4 个 DMA 总线通道组，分别为 DMA\_0 通道组、DMA\_1 通道组、DMA\_2 通道组和 DMA\_3 通道组。其中 DMA\_0 通道组和 DMA\_1 通道组为保留通道组，不可使用；DMA\_2 通道组的 DMA 总线是 AXI 矩阵访问 HP L2MEM 的总线，DMA\_3 通道组的 DMA 总线是 AHB 矩阵访问 HP L2MEM 的总线，详情请见章节 \ref{mod:sysmem} \textit{\nameref{mod:sysmem}} > \ref{other:sysmem-internal-sram-0} 小节。
DMA 总线访问记录（L2MEM 监控功能）的配置流程如下：


\begin{enumerate}

    \item 配置监测地址范围
        \begin{itemize}
            \item 配置 \hyperref[regdesc:L2MEMMONITORLOGMINREG]{L2\_MEM\_MONITOR\_LOG\_MIN\_REG} 和 \hyperref[regdesc:L2MEMMONITORLOGMAXREG]{L2\_MEM\_MONITOR\_LOG\_MAX\_REG} 确定地址范围。 监测地址范围需处于 0x4FF0\_0000 \verb+~+ 0x4FFB\_FFFF 或 0x8FF0\_0000 - 0x8FFB\_FFFF 区间内。
        \end{itemize}
    \item 配置监测模式 (\hyperref[fielddesc:L2MEMMONITORLOGMODE]{L2\_MEM\_MONITOR\_LOG\_MODE})
       \begin{itemize}
            \item 写监测（监测总线是否发生写操作）
           \item 字 (word) 监测（监测总线是否写某个特定 word）
           \item 半字 (halfword) 监测（监测总线是否写某个特定 halfword）
           \item 字节 (byte) 监测（监测总线是否写某个特定 byte）
       \end{itemize}
    \item 配置需要监测的特殊值
        \begin{itemize}
            \item 监测模式为 word 监测时，\hyperref[regdesc:L2MEMMONITORLOGCHECKDATAREG]{L2\_MEM\_MONITOR\_LOG\_CHECK\_DATA\_REG} 即为监测的特定 word。
            \item 监测模式为 halfword 监测时，\hyperref[regdesc:L2MEMMONITORLOGCHECKDATAREG]{L2\_MEM\_MONITOR\_LOG\_CHECK\_DATA\_REG} 的 [15:0] 即为监测的特定 halfword。
            \item 监测模式为 byte 监测时，\hyperref[regdesc:L2MEMMONITORLOGCHECKDATAREG]{L2\_MEM\_MONITOR\_LOG\_CHECK\_DATA\_REG} 的 [7:0] 即为监测的特定 byte。
            \item \hyperref[regdesc:L2MEMMONITORLOGDATAMASKREG]{L2\_MEM\_MONITOR\_LOG\_DATA\_MASK\_REG} 用于屏蔽 \hyperref[regdesc:L2MEMMONITORLOGCHECKDATAREG]{L2\_MEM\_MONITOR\_LOG\_CHECK\_DATA\_REG} 的相应 byte。当某一 byte 被屏蔽，表示该 byte 可为任意的值，比如 word 监测，\hyperref[regdesc:L2MEMMONITORLOGCHECKDATAREG]{L2\_MEM\_MONITOR\_LOG\_CHECK\_DATA\_REG} 配置为 0x01020304，\hyperref[regdesc:L2MEMMONITORLOGDATAMASKREG]{L2\_MEM\_MONITOR\_LOG\_DATA\_MASK\_REG} 配置为 0x1，因此只要总线写 0x010203XX 都会被记录。
        \end{itemize}
    \item 配置记录信息的存储地址范围
        \begin{itemize}
            \item \hyperref[regdesc:L2MEMMONITORLOGMEMSTARTREG]{L2\_MEM\_MONITOR\_LOG\_MEM\_START\_REG} 和 \hyperref[regdesc:L2MEMMONITORLOGMEMENDREG]{L2\_MEM\_MONITOR\_LOG\_MEM\_END\_REG} 为存储地址范围配置寄存器。 记录信息的存储地址范围为 0x4FF0\_0000 \verb+~+ 0x4FFB\_FFFF。
            \item 置位 \hyperref[regdesc:L2MEMMONITORLOGMEMADDRUPDATEREG]{L2\_MEM\_MONITOR\_LOG\_MEM\_ADDR\_UPDATE\_REG}，使 \hyperref[regdesc:L2MEMMONITORLOGMEMCURRENTADDRREG]{L2\_MEM\_MONITOR\_LOG\_MEM\_\\CURRENT\_ADDR\_REG} 更新为 \hyperref[regdesc:L2MEMMONITORLOGMEMSTARTREG]{L2\_MEM\_MONITOR\_LOG\_MEM\_START\_REG}。
            \item 配置辅助调试模块的 L2MEM 监控功能对 HP L2MEM 的使用权限，只有开启了辅助调试模块对 HP L2MEM 的使用权限才能访问 HP L2MEM。关于如何配置，详情见章节 \ref{mod:permctrl} \textit{\nameref{mod:permctrl}} 的 \hyperref[regdesc:TEEDMAL2MEMMONPMSWREG]{TEE\_DMA\_L2MEM\_MON\_PMS\_W\_REG} 及相关内容。
        \end{itemize}
    \item 配置存储记录信息的内存的存储模式：loop 模式和非 loop 模式
        \begin{itemize}
            \item loop 模式下，将记录信息循环写进配置地址范围内，当写到结束地址时回到起始地址开始写，覆盖之前记录的信息。置位 \hyperref[fielddesc:L2MEMMONITORLOGMEMLOOPENABLE]{L2\_MEM\_MONITOR\_LOG\_MEM\_LOOP\_ENABLE} 使用 loop 模式。\\
            \item 非 loop 模式下，当将记录信息写到结束地址后，会一直停留在结束地址，并不再写记录信息，不会覆盖最开始的记录信息，并且丢弃后面的记录信息。清零 \hyperref[fielddesc:L2MEMMONITORLOGMEMLOOPENABLE]{L2\_MEM\_MONITOR\_LOG\_MEM\_LOOP\_ENABLE} 使用非 loop 模式。\\
            \item 示例同章节 \ref{sec:hp_cpu_access_mon} 的步骤 \ref{other:dbgassist-item-log-config-step5}。
        \end{itemize}
     \item 配置总线使能
        \begin{itemize}
            \item 配置 \hyperref[fielddesc:L2MEMMONITORLOGDMA2ENA]{L2\_MEM\_MONITOR\_LOG\_DMA\_2\_ENA} 和 \hyperref[fielddesc:L2MEMMONITORLOGDMA3ENA]{L2\_MEM\_MONITOR\_LOG\_DMA\_3\_ENA} 分别使能 DMA\_2 通道组和 DMA\_3 通道组总线访问记录，可同时使能。
        \end{itemize}
\end{enumerate}

辅助调试模块会先将 DMA 记录信息缓存到一个内部 buffer B 里，然后从 buffer 取出记录信息再存储到配置的存储区域内。当监测行为被连续触发时，连续产生的记录数据包可能会使得 buffer 满而不能缓存新的记录数据包，此时模块会舍弃掉不能及时缓存到 buffer 的记录数据包，并在 buffer 不满的时候缓存一个 DMA LOST 数据包作为替代，此时仅能知道 DMA LOST 数据包生成时所舍弃的记录数据包的总线类型，但不能知道在此之前舍弃了多少个记录数据包。


当总线访问记录结束后需要从内存中读取记录数据并进行解码。记录数据里只有三种包格式，DMA\_2 数据包、DMA\_3 数据包和 DMA LOST 数据包，DMA\_2 数据包对应 DMA\_2 通道组数据总线记录信息，DMA\_3 数据包对应 DMA\_3 通道组数据总线记录信息，三种包格式如表 \ref{tab:dbgassist-dma2-package}、\ref{tab:dbgassist-dma3-package} 和 \ref{tab:dbgassist-dma-lost-package} 所示：

    \begin{longtable}[c]{|L{3cm}|L{2cm}|L{2cm}|l|}
    \caption{DMA\_2 包格式}\label{tab:dbgassist-dma2-package}
    \\\hline
    \rowcolor{lightgray}
    {\bfseries Bit[31:9]} & {\bfseries Bit[8:4]} & {\bfseries Bit[3:1]} & {\bfseries Bit[0]} \\\hline
     addr\_offset0 & reserved & format & anchored(0)  \\\hline
    \end{longtable}

    \begin{longtable}[c]{|L{3cm}|L{2cm}|L{2cm}|l|}
    \caption{DMA\_3 包格式}\label{tab:dbgassist-dma3-package}
    \\\hline
    \rowcolor{lightgray}
    {\bfseries Bit[31:9]} & {\bfseries Bit[8:4]} & {\bfseries Bit[3:1]} & {\bfseries Bit[0]} \\\hline
     addr\_offset1 & reserved & format & anchored(0)  \\\hline
    \end{longtable}

    \begin{longtable}[c]{|L{3cm}|L{3cm}|L{3cm}|L{3cm}|l|}
    \caption{DMA LOST 包格式}\label{tab:dbgassist-dma-lost-package}
    \\\hline
    \rowcolor{lightgray}
    {\bfseries Bit[31:10]} & {\bfseries Bit[9:8]} &{\bfseries Bit[7:4]} & {\bfseries Bit[3:1]} & {\bfseries Bit[0]} \\\hline
    reserved & lost\_source & reserved & format & anchored(0) \\\hline
    \end{longtable}

从包格式可以看出，DMA\_2 、DMA\_3 和 DMA LOST 数据包均为 32 位。下面介绍各个域的含义：
\begin{itemize}
    \item {\bfseries format} 表示此次数据包的类型，4 表示 DMA\_2 数据包，5 表示 DMA\_3 数据包，6 表示 DMA LOST 数据包，其他值为保留值。
    \item {\bfseries addr\_offset0} 记录了此次写操作的地址偏移，实际地址 = addr\_offset0 + \{\hyperref[regdesc:L2MEMMONITORLOGMINREG]{L2\_MEM\_MONITOR\_LOG\_MIN\_REG}[31:5]，5'h0\}。
    \item {\bfseries addr\_offset1} 记录了此次写操作的地址偏移，实际地址 = addr\_offset1 + \{\hyperref[regdesc:L2MEMMONITORLOGMINREG]{L2\_MEM\_MONITOR\_LOG\_MIN\_REG}[31:2]，2'h0\}。
    \item {\bfseries lost\_source} 记录了 DMA LOST 数据包生成时舍弃了哪些数据包：
    \begin{itemize}
        \item Bit[8]：0 表示未舍弃 DMA\_2 数据包， 1 表示舍弃了 DMA\_2 数据包。
        \item Bit[9]：0 表示未舍弃 DMA\_3 数据包， 1 表示舍弃了 DMA\_3 数据包。
    \end{itemize}
    \item {\bfseries anchored} 表示该 32 位数据在数据包中的位置，0 表示该 32 位数据是数据包的低 32 位。
\end{itemize}

模块内部 buffer 的数据宽度是 32 位。当同时使能 DMA\_2 通道组或 DMA\_3 通道组总线访问记录并且在某一时刻同时生成了记录数据，先将 DMA\_2 的数据包缓存到 buffer 中，再缓存 DMA\_3 的数据包。辅助调试模块会自动地将缓存数据从 buffer 取出，并将取出的数据再以 32 位数据宽度存入指定范围的 memory 里。

数据包解析流程如下：

\begin{itemize}
    \item \hyperref[fielddesc:L2MEMMONITORLOGMEMFULLFLAG]{L2\_MEM\_MONITOR\_LOG\_MEM\_FULL\_FLAG} 确定数据是否溢出配置范围。
    \begin{itemize}
        \item 如果未溢出，则数据读取的地址范围为 \hyperref[regdesc:L2MEMMONITORLOGMEMSTARTREG]{L2\_MEM\_MONITOR\_LOG\_MEM\_START\_REG}  \verb+~+ \hyperref[regdesc:L2MEMMONITORLOGMEMCURRENTADDRREG]{L2\_MEM\_MONITOR\_LOG\_MEM\_CURRENT\_ADDR\_REG} - 4 ；
        \item 若溢出并且开启了 loop 模式，则数据读取的地址范围为 \hyperref[regdesc:L2MEMMONITORLOGMEMCURRENTADDRREG]{L2\_MEM\_MONITOR\_LOG\_MEM\_CURRENT\_ADDR\_REG} \verb+~+ \hyperref[regdesc:L2MEMMONITORLOGMEMENDREG]{L2\_MEM\_MONITOR\_LOG\_MEM\_END\_REG} 和 \hyperref[regdesc:L2MEMMONITORLOGMEMSTARTREG]{L2\_MEM\_MONITOR\_LOG\_MEM\_START\_REG} \verb+~+ \hyperref[regdesc:L2MEMMONITORLOGMEMCURRENTADDRREG]{L2\_MEM\_MONITOR\_LOG\_MEM\_CURRENT\_ADDR\_REG} - 4；
        \item 若溢出并且未开启 loop 模式，则数据读取的地址范围为 \hyperref[regdesc:L2MEMMONITORLOGMEMSTARTREG]{L2\_MEM\_MONITOR\_LOG\_MEM\_START\_REG} \verb+~+ \hyperref[regdesc:L2MEMMONITORLOGMEMENDREG]{L2\_MEM\_MONITOR\_LOG\_MEM\_END\_REG}。
    \end{itemize}
    \item 从起始地址处开始读取数据并解析，每次读取 32 位。

\end{itemize}

数据解析完成之后需要通过配置 \hyperref[fielddesc:L2MEMMONITORCLRLOGMEMFULLFLAG]{L2\_MEM\_MONITOR\_CLR\_LOG\_MEM\_FULL\_FLAG} 清除  \hyperref[fielddesc:L2MEMMONITORLOGMEMFULLFLAG]{L2\_MEM\_MONITOR\_LOG\_MEM\_FULL\_FLAG} 标志位。


% >>> If your module has no registers, delete sample text
%       for Sections Register Summary and Registers below <<<

%%%%%%%%%%%%%%%%%%%%%%%%
%%% Register Summary %%%
%%%%%%%%%%%%%%%%%%%%%%%%

% >>> Update the sample text below and insert
%     the info on register summary into the subfile <<<

\newpage
\hypertarget{dbgassist-reg-summ}{}
\section{寄存器列表}
\label{sec:dbgassist-reg-summ}

本小节中\textbf{HP CPU 总线记录配置寄存器} 的地址为相对于 \textbf{SPM 监控} 基地址的地址偏移量（相对地址），详见 \ref{sec:dbgassist-spm-mem-monitor-reg-summ}，\textbf{DMA 总线记录配置寄存器} 的地址为相对于 \textbf{L2MEM 监控} 基地址的地址偏移量（相对地址），详见 \ref{sec:dbgassist-l2-mem-monitor-reg-summ}，其它寄存器的所有地址均为相对于 \textbf{总线监控} 基地址的地址偏移量（相对地址），详见 \ref{sec:dbgassist-bus-monitor-reg-summ}。具体基地址请见章节 \ref{mod:sysmem} \textit{\nameref{mod:sysmem}} 中的表 \ref{tab:sysmem-base-address}。

请查看章节 \hyperref[glossary-access-types]{\textit{寄存器的访问类型}}，了解“访问”列缩写的含义。

\subsection{HP CPU 总线记录配置寄存器列表}
\label{sec:dbgassist-spm-mem-monitor-reg-summ}
\subfile{\modulefiles/23-DBGASSIST-spm-mem-monitor-reg-summ__CN}

\subsection{DMA 总线记录配置寄存器列表}
\label{sec:dbgassist-l2-mem-monitor-reg-summ}
\subfile{\modulefiles/23-DBGASSIST-l2-mem-monitor-reg-summ__CN}

\subsection{其它寄存器列表}
\label{sec:dbgassist-bus-monitor-reg-summ}


\subfile{\modulefiles/23-DBGASSIST-reg-summ__CN}


%%%%%%%%%%%%%%%%%
%%% Registers %%%
%%%%%%%%%%%%%%%%%

% >>> Update the sample text below and insert
%      the info on registers into the subfile <<<

\newpage
\section{寄存器}


本小节中\textbf{HP CPU 总线记录配置寄存器} 的地址为相对于 \textbf{SPM 监控} 基地址的地址偏移量（相对地址），详见 \ref{sec:dbgassist-spm-mem-monitor-reg}，\textbf{ DMA 总线记录配置寄存器} 的地址为相对于 \textbf{L2MEM 监控} 基地址的地址偏移量（相对地址），详见 \ref{sec:dbgassist-l2-mem-monitor-reg}，其它寄存器的所有地址均为相对于 \textbf{总线监控} 基地址的地址偏移量（相对地址），详见 \ref{sec:dbgassist-bus-monitor-reg}，具体基地址请见章节 \ref{mod:sysmem} \textit{\nameref{mod:sysmem}} 中的表 \ref{tab:sysmem-base-address}。

关于保留 (reserved) 域的处理，请查看章节 \hyperref[programming-reserved-field]{\textit{如何配置寄存器的保留域}}。


\subsection{HP CPU 总线记录配置寄存器}
\label{sec:dbgassist-spm-mem-monitor-reg}

\subfile{\modulefiles/23-DBGASSIST-spm-mem-monitor-reg__CN}


\subsection{DMA 总线记录配置寄存器}
\label{sec:dbgassist-l2-mem-monitor-reg}

\subfile{\modulefiles/23-DBGASSIST-l2-mem-monitor-reg__CN}


\subsection{其它寄存器}
\label{sec:dbgassist-bus-monitor-reg}

\subfile{\modulefiles/23-DBGASSIST-reg__CN}


\end{document}
